\section{Introduction}
Social media platforms face mounting regulatory pressure worldwide. Australia and Spain are implementing age restrictions for users under 16\footnote{Australia approves social media ban on under-16 \url{https://www.bbc.com/news/articles/c89vjj0lxx9o}}, while several countries have banned TikTok due to concerns about misinformation \cite{basch2021global} and harmful content \cite{weimann2023research}. The EU's Digital Service Act mandates that large platforms "identify, analyze and assess any systemic risks [...] from the use made of their services," including "actual or foreseeable negative effects on civic discourse."

Meeting these regulatory requirements demands robust quantitative evidence about platform risks—particularly the role of algorithmic mechanisms like ranking and recommendation systems. Yet obtaining such evidence remains challenging in computational social and communication science \cite{wang2016understanding}. The fundamental problem is causal inference: determining whether alternative platform designs would produce different, potentially more favorable outcomes. This question is best addressed through simulation, which enables testing counterfactual scenarios without exposing real users to potentially harmful conditions.

The validity of simulation-based insights depends critically on how faithfully simulations represent reality. Recent advances in Large Language Models (LLMs) have enabled more sophisticated agent communication in social simulations \cite{chuang2024simulating}, surpassing the realism of traditional symbolic agent languages. However, this increased sophistication raises a crucial question: Can LLM-based agents replicate human social behavior with sufficient fidelity to produce scientifically robust insights?

Answering this question requires systematic experimental designs that improve reproducibility and provide quantitative measures of \textit{empirical realism}, the degree to which simulated behaviors match observed human behaviors. Without such validation, simulation findings lack the credibility needed to inform policy decisions or scientific understanding.

\subsection*{Contributions}
Our work advances the rigor of LLM-based social network simulation through three key contributions:

\paragraph{Framework}
We formalize social networks to (i) standardize their simulation and (ii) enable systematic benchmarking of empirical realism through quantifiable metrics.

\paragraph{Case Study}
We instantiate this framework with multiple approaches for imitating user communication on $\mathbb{X}$, including posting, replying, and reply-likelihood prediction tasks across English and German datasets.

\paragraph{Validation}
We benchmark empirical realism across tasks, languages, and modeling approaches, revealing when and how LLMs can faithfully replicate user behavior, and crucially, when they cannot.


